% SPDX-License-Identifier: CC-BY-4.0
\section{Formal-verification map and reproduction}\label{sec:verification}

The arguments are complete in the body; reading them requires no Lean. This
appendix locates the formal declarations and gives commands to check them. File
paths are relative to \texttt{formalization/}. The table uses the prefix
\texttt{MR.} for \nolinkurl{MathResearch.}; the external inputs lie under
\nolinkurl{MathResearch.ThirdParty.}

\begingroup
\small
\renewcommand{\arraystretch}{1.3}
\begin{longtable}{@{}>{\raggedright\arraybackslash}p{0.30\textwidth}>{\raggedright\arraybackslash}p{0.65\textwidth}@{}}
\toprule
Paper statement & Lean file and declarations\\
\midrule
\endfirsthead
\toprule
Paper statement & Lean file and declarations\\
\midrule
\endhead
\bottomrule
\endfoot
Hasse multiplicity, Definition~\ref{def:mult}; Lemma~\ref{lem:mult-basic}
& \leanref{claims/HasseMultiplicity.lean}{\nolinkurl{claims/HasseMultiplicity.lean}}\newline
\nolinkurl{MR.HasseMultiplicity.mult}, \nolinkurl{coe_le_mult_iff}, \nolinkurl{min_le_mult_add},
\nolinkurl{add_le_mult_mul}, \nolinkurl{mult_mul}, \nolinkurl{mult_eq_zero_iff}, \nolinkurl{mult_rename_equiv}\\
Multiplicity Schwartz--Zippel \cite{DKSS13}, Theorem~\ref{thm:dkss}
& \leanref{third-party-claims/MultiplicitySchwartzZippel.lean}{\nolinkurl{third-party-claims/MultiplicitySchwartzZippel.lean}}\newline
\nolinkurl{MR.ThirdParty.MultiplicitySchwartzZippel.multiplicity_schwartz_zippel}\\
Schwartz--Zippel bound over \(\Ft\), Proposition~\ref{prop:cover-bound}
& \leanref{claims/BinaryMultiplicityCoverBound.lean}{\nolinkurl{claims/BinaryMultiplicityCoverBound.lean}}\newline
\nolinkurl{MR.BinaryMultiplicityCoverBound.cover_bound}, \nolinkurl{cover_bound_of_origin}\\
Linear substitutions, Lemma~\ref{lem:subst}; restriction, Lemma~\ref{lem:restrict}
& \leanref{claims/OneDimensionStep.lean}{\nolinkurl{claims/OneDimensionStep.lean}}\newline
\nolinkurl{MR.OneDimensionStep.homogeneousComponent_subst}, \nolinkurl{totalDegree_subst_le},
\nolinkurl{subst_comp}, \nolinkurl{mult_zero_subst}, \nolinkurl{mult_subst}, \nolinkurl{mult_restrict}\\
Avoiding a linear form, Lemma~\ref{lem:linear-form}
& \leanref{claims/LinearFormAvoidance.lean}{\nolinkurl{claims/LinearFormAvoidance.lean}}\newline
\nolinkurl{MR.LinearFormAvoidance.exists_linForm_not_dvd}, \nolinkurl{exists_subst_linForm_eq_X}\\
One-dimension step, Theorem~\ref{thm:step}
& \leanref{claims/OneDimensionStep.lean}{\nolinkurl{claims/OneDimensionStep.lean}}\newline
\nolinkurl{MR.OneDimensionStep.one_dimension_step}\\
Lower bound, Theorem~\ref{thm:lower}
& \leanref{claims/BinaryMultiplicityLowerBound.lean}{\nolinkurl{claims/BinaryMultiplicityLowerBound.lean}}\newline
\nolinkurl{MR.BinaryMultiplicityLowerBound.lower_bound}\\
Digit sums, Lemma~\ref{lem:digits}
& \leanref{claims/TruncatedProduct.lean}{\nolinkurl{claims/TruncatedProduct.lean}}\newline
\nolinkurl{MR.TruncatedProduct.s2_add_le}, \nolinkurl{s2_two_pow}, \nolinkurl{s2_le_self},
\nolinkurl{two_mul_add_s2_le}; \nolinkurl{MR.BinaryMultiplicityDegreeFormula.s2_two_pow_sub_one}\\
Telescoping, Lemma~\ref{lem:telescope}; truncated product, Theorem~\ref{thm:truncated-product}
& \leanref{claims/TruncatedProduct.lean}{\nolinkurl{claims/TruncatedProduct.lean}}\newline
\nolinkurl{MR.TruncatedProduct.telescope}, \nolinkurl{g}, \nolinkurl{truncated_product}\\
Per-order construction, Theorem~\ref{thm:per-order}
& \leanref{claims/PerOrderConstruction.lean}{\nolinkurl{claims/PerOrderConstruction.lean}}\newline
\nolinkurl{MR.PerOrderConstruction.per_order_construction}\\
Formula below \(2^{n-1}\), Theorem~\ref{thm:main-formal}; upper bound of Theorem~\ref{thm:main}; Definition~\ref{def:admissible}
& \leanref{claims/BinaryMultiplicityDegreeFormula.lean}{\nolinkurl{claims/BinaryMultiplicityDegreeFormula.lean}}\newline
\nolinkurl{MR.BinaryMultiplicityDegreeFormula.Admissible}, \nolinkurl{degree_formula},
\nolinkurl{degree_formula_exact}, \nolinkurl{degree_formula_of_lt_two_pow}\\
Extended range, Corollary~\ref{cor:extended}
& \leanref{claims/BinaryMultiplicityDegreeExtendedRange.lean}{\nolinkurl{claims/BinaryMultiplicityDegreeExtendedRange.lean}}\newline
\nolinkurl{MR.BinaryMultiplicityDegreeExtendedRange.degree_formula_extended}\\
Hyperplane covers \cite{BBDM23}, Theorem~\ref{thm:hyperplane}
& \leanref{third-party-claims/HyperplaneCoverUpperBound.lean}{\nolinkurl{third-party-claims/HyperplaneCoverUpperBound.lean}}\newline
\nolinkurl{MR.ThirdParty.HyperplaneCoverUpperBound.hyperplane_cover_upper_bound},
\nolinkurl{card_dot_eq_one}\\
Main theorem, Theorem~\ref{thm:main}
& \leanref{claims/BinaryMultiplicityDegreeComplete.lean}{\nolinkurl{claims/BinaryMultiplicityDegreeComplete.lean}}\newline
\nolinkurl{MR.BinaryMultiplicityDegreeComplete.degreeMin}, \nolinkurl{degree_complete},
\nolinkurl{cover_bound_attained_iff}\\
\end{longtable}
\endgroup

\paragraph{Reading the formal statements.} In Lean, points are functions
\(\mathrm{Fin}\,n\to\mathrm{ZMod}\,2\), degree is \texttt{MvPolynomial.totalDegree},
and multiplicity is \nolinkurl{MR.HasseMultiplicity.mult}, the order of
\(P(a+z)\) in \(\mathbb N_\infty\). The one-dimension step is stated for the
variables \(\mathrm{Option}\,\sigma\), with \(\mathrm{none}\) playing the role of
\(x_0\). The lower bound relabels \(\mathrm{Fin}(n+1)\) as
\(\mathrm{Option}(\mathrm{Fin}\,n)\).

\paragraph{Commands.} The project pins Lean \texttt{4.34.0-rc2} and Mathlib in
\texttt{lakefile.lean} and \texttt{lake-manifest.json}. With
\href{https://lean-lang.org/install/manual/}{elan}, Git and Python~3 installed:

\begingroup\small
\begin{verbatim}
git clone https://github.com/kbr-/math-research.git bmd-check
cd bmd-check
git checkout --detach 8651228097422970d147972353bbf9a1c6ed6938
cd formalization
export MATHLIB_NO_CACHE_ON_UPDATE=1
lake exe cache get claims/BinaryMultiplicityDegreeComplete.lean
python3 verify.py --target claims/BinaryMultiplicityDegreeComplete.lean \
  --out verification.txt
\end{verbatim}
\endgroup

The target imports every module in the table, so the build covers the whole
chain, and its axiom report is transitive. The verifier builds the proofs, prints each listed
declaration's type, and rejects unfinished proofs and nonstandard axioms. The
expected axiom report for the main theorem is
\begin{quote}\small\ttfamily
'MathResearch.BinaryMultiplicityDegreeComplete.\par
degree\_complete' depends on axioms:\par
[propext, Classical.choice, Quot.sound]
\end{quote}
These are Lean's standard foundations.
